\section{Conclusion and Future Work}
\label{sec:conclusion}

We have presented \alaya{}, a memory engine for conversational AI agents
that treats memory as a dynamic process rather than a static store. Drawing
on complementary learning systems theory, Bjork dual-strength forgetting,
Hebbian learning, and Yogac\=ara Buddhist psychology, \alaya{} implements
seven lifecycle operations---consolidation, transformation, forgetting,
preference emergence, emergent ontology, retrieval-induced forgetting, and
cascade purge---that manage memory health over extended interaction
horizons.

The emergent ontology contribution is particularly novel: categories arise
from dual-signal clustering (embedding similarity + graph topology) during
the transformation lifecycle, without supervision or manual curation.
In v0.2.0, categories support hierarchical nesting: when a category grows
too large and internally incoherent, it automatically splits into
sub-categories, producing an is-a taxonomy that refines as knowledge
grows. Categories participate in the Hebbian graph via \texttt{MemberOf}
links, enabling cross-domain bridging where spreading activation flows
through category nodes to surface categorically related memories during
retrieval.

Our baseline evaluation reveals a statistically significant retrieval
crossover between full-context injection and na\"ive RAG---neither extreme
is universally optimal---motivating the lifecycle-aware memory management
that \alaya{} provides. The system is implemented as a zero-dependency Rust
library requiring no external databases or network services, available as
an open-source crate at \url{https://crates.io/crates/alaya}.

\subsection*{Future Work}

\textbf{Prototype theory and category seeds.} The current centroid-based
category membership treats all members equally. Future work will add
graded typicality (\emph{n\=ama-r\=upa}---name and form), where each
member has a typicality score reflecting how representative it is of its
category. Category seeds (\emph{b\=ija}) will retain dormant state,
allowing dissolved categories to re-emerge when new evidence arrives
rather than being discovered from scratch.

\textbf{End-to-end evaluation} of \alaya{}'s full retrieval pipeline
(BM25 + vector + graph + Hebbian reranking) with lifecycle operations
(consolidation, forgetting, category-boosted retrieval) against LoCoMo
and LongMemEval baselines is the immediate empirical priority.

\textbf{Learned thresholds} for category assignment, discovery, merge,
splitting, and dissolution---replacing the current hand-tuned defaults
with values optimized from task performance---represent a natural
extension of the RL-trained memory management trend identified in the
companion survey~\cite{hui2026taxonomy}.
